\documentclass[11pt]{article}
\usepackage[a4paper, top=18mm, bottom=18mm, left=20mm, right=20mm]{geometry}
\usepackage[T2A]{fontenc}
\usepackage[utf8]{inputenc}
\usepackage[ukrainian]{babel}
\usepackage{graphicx}
\usepackage[export]{adjustbox}
\usepackage{amsmath}
\usepackage{amssymb}
\graphicspath{{/app/Quiz/Scripts/2023/ProgAll/15BinTree/}{/app/Quiz/Scripts/2021/Mod2021_3/BinTree/}{/app/Quiz/Scripts/2020/Mod2020_4/BinTree/}}
\setlength{\parindent}{0pt}
\setlength{\parskip}{6pt}
\pagestyle{empty}
\begin{document}
%begin
{\large\textbf{Контрольна робота}}

\textbf{Варіант \No\,4}\hfill \textit{ПІБ:}\,\underline{\hspace{60mm}}
\vspace{4mm}

%beginitem
1. 

print('R', end=' ')
for b in range(-3,-3,3):
    if b<=4:
        continue
        print(b, end=' ')
        b=2
    if b<=4:
        break
    else:
        print(b, end=' ')
    print(b)
\vspace{5mm}
%enditem

%beginitem
2. 
a = 1
b = 4
c = 7

def g(b):
global c    a -= 5
    b = 1
    c = 3
    return a + b + c

a = 0
b = 9
c = 5

try:
    print(g(a), a, b, c, sep=';')
except:
    print('Z', a, b, c, sep=';')
\vspace{5mm}
%enditem

%beginitem
3. 

\begin{center}
	\adjustimage{max size={0.8\linewidth}{0.8\paperheight}}
	{../15BinTree/trees_exp/35.png}
\end{center}
Указати послідовність міток вершин дерева, зображеного на рис., що утвориться в результаті обходу дерева в оберненому порядку. Мітки записувати через пробіл.\par Алгоритм починає роботу з кореня (на рис. це верхня вершина).
%answer:35 оберненому
\vspace{5mm}
%enditem

%beginitem
4. 
Записати ВИРАЗ, значенням якого є список, що складається з квадратних коренівцілих чисел від 14 до 2012 (включно), що діляться на 3, записаних в оберненому порядку.\vspace{5mm}
%enditem

%beginitem
5. %goodpath:Scripts/2019/Mod2019_1/Lex/01-good.txt
Відмітити все, що є однією правильною лексемою:
( 1 ) 'c'

( 2 ) 0123

( 3 ) 0b23

( 4 ) xy
\vspace{5mm}
%enditem

%beginitem
6. 
Написати програму, що запитує у користувача значення дійсних параметрів $x$ та $y$, обчислює для них значення виразу 
$$\frac{\pi x - 19}{(\arcsin x - 0.5) (y - 23)}$$ 
та повiдомляє введенi данi та результат обчислення.
 Оцінюється правильність та якість коду.

\vspace{5mm}
%enditem

%beginitem
7. 
b=['0','1','2','3','4']; b.append('13'); print(b)\vspace{5mm}
%enditem

%beginitem
8. 
Записати ВИРАЗ, значенням якого є множина, що складається з цілих чисел від 14 до 2012 (включно), що діляться на 3.\vspace{5mm}
%enditem

%beginitem
9. %goodpath:Scripts/2018/Mod2018_1/01-good.txt
Відмітити все, що є однією правильною лексемою:
( 1 ) a2f

( 2 ) 'qwerty'

( 3 ) **

( 4 ) true
\vspace{5mm}
%enditem

%beginitem
10. %goodpath:Scripts/2019/Mod2019_1/Lex/03-good.txt
Відмітити все, що є коректним оператором С++:
( 1 ) return 'a'<'c'

( 2 ) print(sep=":","qwe")

( 3 ) x=*2

( 4 ) float(2)/87
\vspace{5mm}
%enditem

%end
\newpage
\end{document}

